% Project journal: Real-Time Bayesian Experimentation Engine
\documentclass[11pt,a4paper]{article}
\usepackage[margin=2.2cm]{geometry}
\usepackage{amsmath,amssymb}
\usepackage{booktabs}
\usepackage{enumitem}
\usepackage{xcolor}
\usepackage{graphicx}
\usepackage[hidelinks]{hyperref}
\setlist{nosep,leftmargin=1.4em}
\setlength{\parindent}{0pt}
\setlength{\parskip}{5pt}

\newcommand{\Rs}{Rs.\,}
\newcommand{\entry}[2]{\section*{#1 \hfill\normalfont\small #2}\addcontentsline{toc}{section}{#1}}
\newcommand{\jpart}[1]{\subsection*{#1}}
\newcommand{\code}[1]{\texttt{#1}}

\title{Project Journal\\[4pt]\large Real-Time Bayesian Experimentation Engine}
\date{7--20 October 2026}
\author{}

\begin{document}
\maketitle

\begin{abstract}
\noindent This journal records, day by day, how a Bayesian experimentation engine for choosing checkout
offers on a quick-commerce app was designed, built and evaluated. It covers the experiment design, the
traffic simulator, the statistics behind each method, the bugs found along the way, and the measured
results of five decision methods across three scenarios. All data is simulated, so every method is scored
against a known truth.
\end{abstract}

\tableofcontents

% ======================================================================
\entry{7 October: The first plan}{planning}

\jpart{Discussion}
The project began as a Bayesian demand-forecasting and inventory plan built on a Poisson--Gamma model and the
Walmart M5 dataset. On review, the plan was set aside for three reasons:
\begin{itemize}
  \item M5 is stored in wide format and is highly intermittent, with many zero-sales days.
  \item Sales are censored demand: a shelf that is empty records zero sales even when customers want the product.
  \item The inventory side would have to be simulated anyway, and the plan assumed more modelling experience
        than the starting point, which was Beta--Binomial updating.
\end{itemize}

% ======================================================================
\entry{15 October: Pivot to experimentation}{planning, research}

\jpart{Discussion}
An earlier project (DataFlow) read as backend engineering rather than data science: it posed no question,
fitted no model and drew its insights from data it had generated itself. The new project is a
\textbf{Bayesian experimentation engine} that chooses which checkout offer to show on a quick-commerce app.
It builds directly on Beta updating, and experimentation is a core skill for data science and product
analytics roles. Target date: 25 October.

\jpart{Research}
\begin{itemize}
  \item Fixed-horizon A/B tests remain the industry default. Checking results repeatedly (``peeking'')
        inflates false positives (Evan Miller).
  \item Industry responses include sequential tests (Spotify), always-valid inference (Optimizely) and
        variance reduction with pre-experiment data (CUPED at Microsoft Bing, about 50\% less variance).
  \item Bayesian A/B testing with an expected-loss stopping rule is used in commercial tools (VWO).
        Thompson sampling (Thompson, 1933; Chapelle and Li at Yahoo, 2011) powers bandit systems at
        Google, Netflix, Stitch Fix and JD.com.
  \item Bandits are good at earning while learning and weaker at measuring effect sizes. A fair study must
        compare them with a proper A/B test, not a weakened one.
\end{itemize}

% ======================================================================
\entry{18 October: Experiment design and the simulator}{design}

\jpart{The experiment}
Which checkout offer should a campaign show? The unit is one cart-page visit; the outcome is whether the
visitor orders and, if so, the profit earned. The toy values below are assumptions, not real data.

\begin{center}
\begin{tabular}{clccc}
\toprule
Offer & Description & True order rate & Profit per order & Profit per visitor \\
\midrule
0 & No offer (control)   & 0.08 & \Rs 90 & \Rs 7.20 \\
1 & Free delivery        & 0.10 & \Rs 70 & \Rs 7.00 \\
2 & \Rs 50 off           & 0.12 & \Rs 40 & \Rs 4.80 \\
3 & \Rs 20 off above 299 & 0.11 & \Rs 70 & \textbf{\Rs 7.70} \\
\bottomrule
\end{tabular}
\end{center}
Offer 2 converts best but earns least, and offer 3 is the true profit winner. A method that judges by
conversion alone picks the wrong offer.

\jpart{Design decisions}
\begin{itemize}
  \item The truth is one order rate per offer, hidden from every decision method and used only for scoring.
  \item Three ideas must be kept apart: the true rate $p$ (fixed and hidden), the outcome of a visit
        (a Bernoulli draw), and the belief about $p$ (a Beta posterior).
  \item Returning users must always see the same offer. Assignment uses a salted hash,
        \\hspace*{1em}\code{sha256("user\_id:experiment\_id") mod 4};\ plain \code{user\_id \% 4} would give two
        experiments identical splits, so their effects could not be separated.
  \item Arrivals follow a non-homogeneous Poisson process: each minute has its own rate $\lambda_t$.
        The daily total is Poisson with mean $\sum_t\lambda_t$, but the timing comes from how $\lambda_t$ varies.
  \item ``Returning-visit share'' (the share of visits from someone already seen in the window) differs
        from the customer-level ``repeat purchase rate''. A share of 0.15 is used as an assumption.
\end{itemize}

\jpart{Calculations}
\begin{align*}
\text{relative spread of a Poisson count} &= \frac{\sqrt{\lambda}}{\lambda} = \frac{1}{\sqrt{\lambda}}
  \qquad (\lambda=40:\ 16\%;\ \lambda=2400:\ 2\%)\\
\text{expected distinct users (pool $M$, $N$ visits)} &= M\left(1-e^{-N/M}\right)\\
\text{share of visits from someone already seen} &= 1-\frac{1}{\text{average visits per active user}}
\end{align*}

\jpart{Daily traffic}
Daypart shares follow ranges from an industry article (Zane, 2 August 2026), which cites no dataset and is
therefore treated as a plausible range: 07--11: 16\%, 11--17: 24\%, 17--23: 46\%, 23--02: 10\%, and an
assumed floor of 4\% for 02--07. Volume is about 20{,}000 visits per day.

% ======================================================================
\entry{19 October: Simulator, first tests and the fixed A/B baseline}{code, results}

\jpart{The per-visit simulator}
\begin{itemize}
  \item Bugs found and fixed: a NumPy shape passed as a data type; \code{rng.shuffle} returning \code{None};
        \code{np.append} not changing the array in place (a silent bug that made the returning share 0\%);
        a hidden cap of 100{,}000 user IDs that would run out on day six; and Poisson counts drawn once and
        copied to every day.
  \item Verified over 3 days (seed 42, 60{,}374 visits): visits per day 20{,}320, 19{,}884 and 20{,}170;
        window shares 4.0, 16.1, 24.0, 46.0 and 9.9\%; returning share 15.1\%; no user saw more than one offer.
\end{itemize}

\jpart{The two-proportion z-test}
\[
\bar p = \frac{x_1+x_2}{n_1+n_2},\qquad
SE = \sqrt{\bar p(1-\bar p)\left(\tfrac1{n_1}+\tfrac1{n_2}\right)},\qquad
z = \frac{\hat p_2-\hat p_1}{SE},\qquad p\text{-value} = 2\,\Phi^{c}(|z|)
\]
\begin{itemize}
  \item The rate is pooled because, under the null hypothesis, both groups share one rate.
  \item $z$ grows with $\sqrt{n}$: ten times the data gives about 3.16 times the $z$.
  \item The survival function \code{norm.sf} keeps precision in the far tail, where \code{1 - norm.cdf}
        rounds to zero.
  \item A p-value is the probability of data this extreme if the null is true, not the probability that
        the null is true.
\end{itemize}

\jpart{Tests on one 3-day run}
\begin{center}
\small
\begin{tabular}{lcl}
\toprule
Comparison & Result & Note \\
\midrule
Conversion, offer 2 vs control & $z=9.35$ & matches statsmodels; $\chi^2=z^2=87.43$ \\
Conversion, offer 3 vs control & $z=7.00$ & 95\% interval (1.71, 3.04) points \\
Profit per visitor, offer 3 vs control & $t=-0.10$, $p=0.92$ & a true gap of \Rs 0.50 is invisible \\
\bottomrule
\end{tabular}
\end{center}

\jpart{Welch's t-test from totals}
Each visitor earns either 0 or the offer's fixed profit $m$, so the test needs only totals:
\[
\text{mean} = m\,p,\qquad \text{variance} = m^2\,p(1-p),\qquad
t = \frac{\bar x_k-\bar x_0}{\sqrt{s_k^2/n_k + s_0^2/n_0}}.
\]
Detecting a \Rs 0.50 gap with 80\% power at the 5\% level needs
$n \approx (1.96+0.84)^2(\sigma_0^2+\sigma_3^2)/0.5^2 \approx 34{,}000$ visitors per offer, about seven
days of all traffic.

\jpart{The fast simulator and policy classes}
Hourly batches replace the per-visit loop: visits $\sim$ Poisson(hourly rate), the split across offers
$\sim$ Multinomial, and orders $\sim$ Binomial. Each decision method is a class with four methods
(\code{allocate}, \code{update}, \code{should\_stop}, \code{choice\_test}), so the simulation loop never changes.
Monte Carlo runs repeat each campaign 500 times with independent seeds from \code{SeedSequence(42).spawn}.

\jpart{Fixed A/B test: results}
\begin{center}
\begin{tabular}{lcc}
\toprule
Fixed A/B test & 3 days & 14 days \\
\midrule
Normal: picks offer 3 & 32.6\% & 95.4\% \\
A/A: declares a non-control winner & 1.4\% & 2.2\% \\
Tiny gap: picks offer 3 & 3.6\% & 16.2\% \\
\bottomrule
\end{tabular}
\end{center}
All results agree with theory: about 30\% and 95\% power in the normal scenario, about 4\% and 15\% in the
tiny-gap scenario, and about 2\% false alarms at either length. The A/A rate follows from the decision rule:
each comparison needs $t>2.39$ (two-sided $0.05/3$, ``better'' direction only), which happens with
probability 0.83\%; three comparisons that share a control group give about 2\%.

\jpart{Peeking}
Peeking means running the test every evening and stopping at the first significant result. Each extra
look is another chance for noise to cross the threshold. Measured over 14 days, false alarms rise from
2.2\% to 9.4\%, while regret falls by 62\% because the large true gap is usually found early.

\jpart{Notes}
\begin{itemize}
  \item Monte Carlo simulation (repeating a random experiment and counting outcomes) is not MCMC
        (Markov chain Monte Carlo, a method for sampling posterior distributions).
  \item Regret is the profit lost compared with always showing the best offer. The best offer maximises
        rate $\times$ profit, not profit alone. In A/A every offer is best, so regret is always zero and
        A/A is judged by false alarms instead.
\end{itemize}

% ======================================================================
\entry{20 October: Bayesian methods and the final comparison}{code, results}

\jpart{Bayesian expected-loss stopping}
Instead of asking whether a gap is significant, this method asks: \emph{if we pick this offer now and we
are wrong, how much profit per visitor do we expect to lose?} Each offer's order rate has the posterior
Beta$(1+\text{orders},\,1+\text{visits}-\text{orders})$. The method draws 5{,}000 possible worlds,
multiplies each rate by its profit per order, measures each offer's loss against the best offer in every
world, and averages. Each evening it stops if the smallest expected loss is below \Rs 0.01 per visitor.
The sampling error of this estimate is about \Rs 0.0003, well below the threshold.

\jpart{Thompson sampling}
Thompson sampling changes the traffic split while the test runs: every hour, each offer receives a share
of traffic equal to its probability of being the best offer, estimated as the share of 200 posterior
worlds it wins. Poor offers lose their traffic within hours, while uncertain offers keep some, so
exploration follows uncertainty. In one campaign, offer 2 received 409 visitors in three days instead of
about 15{,}000. The final method combines Thompson's traffic split with the expected-loss stopping rule.

\jpart{Bugs fixed}
Beta parameters fixed at (1,1) instead of updated from data; comparing order rates instead of profit;
averaging the whole loss table instead of each column; an attribute that shadowed a method of the same
name; comparing arrays inside \code{if}; storing a reference to a mutable array instead of its total; and
global variables that leaked from one campaign to the next.

\jpart{Final results (500 campaigns per cell)}
\begin{center}
\small
\begin{tabular}{lccccc}
\toprule
Normal scenario & Fixed & Peeking & Bayesian & Thompson & Thompson + loss \\
\midrule
Picks offer 3, 3 days       & 32.6\% & 39.0\% & 97.2\% & 98.0\% & \textbf{98.6\%} \\
Picks offer 3, 14 days      & 95.4\% & 96.6\% & 99.6\% & \textbf{100\%} & 99.2\% \\
Regret, 3 days (\Rs)        & 61{,}524 & 54{,}631 & 47{,}473 & 8{,}738 & \textbf{8{,}355} \\
Regret, 14 days (\Rs)       & 287{,}006 & 110{,}267 & 62{,}144 & 13{,}430 & \textbf{10{,}197} \\
Profit per visitor, 14 days & 6.68 & 7.31 & 7.48 & 7.65 & \textbf{7.66} \\
Conversion, 14 days         & 10.3\% & 10.7\% & 10.8\% & 10.8\% & 10.9\% \\
\midrule
A/A: non-control pick, 14 days & 2.2\% & 9.4\% & 73.4\% & 74.8\% & 75.2\% \\
Tiny gap: picks offer 3, 14 days & 16.2\% & 26.2\% & 77.0\% & 85.8\% & \textbf{86.0\%} \\
Tiny gap: regret, 14 days (\Rs)  & 37{,}801 & 33{,}747 & 31{,}672 & 23{,}138 & \textbf{22{,}670} \\
\bottomrule
\end{tabular}
\end{center}

\jpart{Interpretation}
\begin{itemize}
  \item The combined engine loses about 28 times less profit than the fixed A/B test over 14 days and earns
        within 0.5\% of the best possible profit per visitor (\Rs 7.70).
  \item In a 3-day campaign, every Bayesian method finds the best offer at least 97\% of the time, against
        33\% for the fixed test.
  \item No method chose offer 2: decisions on profit avoid the conversion trap.
  \item In A/A, the Bayesian methods pick among identical offers at random, at no cost, and stop early in
        only 21--24\% of campaigns (median day 8). They are decision tools, not tools for certifying an effect.
  \item With a 2.5\% gap, the fixed test finds the better offer in 16\% of 14-day campaigns; Thompson-based
        methods find it in 86\%.
\end{itemize}

\jpart{Speed}
Per 14-day campaign: fixed 5\,ms, peeking 4\,ms, Bayesian 7\,ms, and Thompson 183\,ms with 1{,}000 worlds
per hour, reduced to 27\,ms by using 200 worlds.

\jpart{Future scope}
\begin{itemize}
  \item Run campaigns in parallel across CPU cores (about ten times faster, identical results).
  \item Simulate all campaigns at once as arrays (50--100 times faster; requires rewriting the loop and policies).
  \item Serve the engine live: each offer's posterior held in Redis, with FastAPI endpoints to assign an
        offer and record an outcome.
\end{itemize}

\end{document}
